\section{Beta integrals, Euler factors, and gamma identities}
\label{sec:beta}

We now evaluate the endpoint-deleted zero mode at the Tonelli destinations
constructed in \Cref{sec:endpoint}.  Put
\begin{equation}\label{eq:lambda-z-qv}
 \lambda=\alpha+u,\qquad z=\lambda+\gamma+s,\qquad
 q_1=1-z,\qquad v_1=1-\lambda+\beta+s.
\end{equation}

For $r>0$, ordinary beta integration gives
\begin{align}
 &\int_0^\infty y^{-1/2-\lambda-it}
 (y+r)^{-1/2-\gamma+it-s}\dd y \notag\\
 &\qquad=r^{-z}
 \frac{\Gamma(\frac12-\lambda-it)\Gamma(z)}
      {\Gamma(\frac12+\gamma-it+s)}.
 \label{eq:beta-positive}
\end{align}
A genuine absolute-convergence domain, including the subsequent $\rho$- and
$m$-sums, is
\begin{equation}\label{eq:beta-positive-domain}
 \Re\lambda<\frac12,\qquad \Re z>1,\qquad \Re v_1>1.
\end{equation}
For $r=-R<0$, substitute $y=R+x$ to obtain
\begin{align}
 &\int_0^\infty(x+R)^{-1/2-\lambda-it}
 x^{-1/2-\gamma+it-s}\dd x \notag\\
 &\qquad=R^{-z}
 \frac{\Gamma(\frac12-\gamma+it-s)\Gamma(z)}
      {\Gamma(\frac12+\lambda+it)}.
 \label{eq:beta-negative}
\end{align}
Here a genuine domain is
\begin{equation}\label{eq:beta-negative-domain}
 \Re(\gamma+s)<\frac12,\qquad \Re z>1,\qquad \Re v_1>1.
\end{equation}
The two signs need not have a common elementary chamber near the intended
shifts; \Cref{sec:endpoint} proves the two identities separately and then
selects their common meromorphic continuation by regulated exact equality.

The sum over $\rho$ supplies $\zeta(z)$.  We use
\begin{equation}\label{eq:chi-def}
 \zeta(z)=\chi(z)\zeta(1-z),\qquad
 \chi(z)=\pi^{z-1/2}\frac{\Gamma((1-z)/2)}{\Gamma(z/2)}.
\end{equation}
The remaining gcd sum is
\begin{equation}\label{eq:C-hk}
                   C_{h,k}(s,u)
                   =\sum_{m\geq1}\frac{(h,km)^{q_1}}{m^{v_1}}.
\end{equation}

\begin{lemma}[Finite Euler factors]\label{lem:finite-euler}
In the common meromorphic continuation,
\begin{equation}\label{eq:finite-euler-identity}
 h^{-1/2+\gamma+s}k^{-1/2+\lambda}
 \zeta(1-z)C_{h,k}(s,u)
 =Z_{-\gamma-s,\,\beta+s,\,-\lambda}(h,k).
\end{equation}
\end{lemma}

\begin{proof}
Write $g=(h,k)$, $h=gh_0$, $k=gk_0$.  Since
$(h,km)=g(h_0,m)$,
\begin{equation}\label{eq:C-euler}
 C_{h,k}=g^{q_1}\zeta(v_1)
 \prod_{p^\nu\parallel h_0}(1-p^{-v_1})
 \sum_{j\geq0}p^{q_1\min(j,\nu)-v_1j}.
\end{equation}
For $X=p^{-q_1}$ and $Y_0=p^{-v_1}$, splitting the sum at $j=\nu$ gives
\begin{equation}\label{eq:local-euler-identity}
 (1-Y_0)\sum_{j\geq0}X^{-\min(j,\nu)}Y_0^j
 =X^{-\nu}(1-X)(1-Y_0)
 \sum_{i+j\geq\nu}X^iY_0^j.
\end{equation}
This is exactly the local factor in \eqref{eq:Z-euler}, after the exterior
powers in \eqref{eq:finite-euler-identity} are restored.
\end{proof}

Combine the two signs in the archimedean factor
\begin{align}
 \Phi_{\lambda;\beta,\gamma}(s,t)
 ={}&\chi(z)\Gamma(z)g_{\beta,\gamma}(s,t)
 \bigg\{
 \frac{\Gamma(\frac12-\lambda-it)}
      {\Gamma(\frac12+\gamma-it+s)} \notag\\
 &\hspace{42mm}+
 \frac{\Gamma(\frac12-\gamma+it-s)}
      {\Gamma(\frac12+\lambda+it)}
 \bigg\}.
 \label{eq:Phi-def}
\end{align}
There is no additional $\pi^{-s}$ in \eqref{eq:Phi-def}, because it is
already part of $g_{\beta,\gamma}$.

\begin{lemma}[Archimedean pole collapse]\label{lem:gamma-collapse}
Let $a_0=\frac12-\lambda-it$.  Then
\begin{align}
 &\chi(z)\Gamma(z)\left\{
 \frac{\Gamma(a_0)}{\Gamma(a_0+z)}
 +\frac{\Gamma(1-a_0-z)}{\Gamma(1-a_0)}\right\}
 \notag\\
 &\qquad=(2\pi)^z\frac{\Gamma(a_0)}{\Gamma(a_0+z)}
 \frac{\sin\pi(a_0+z/2)}{\sin\pi(a_0+z)}.
 \label{eq:gamma-collapse}
\end{align}
For $t\asymp T$ and $|\Im\lambda|+|\Im s|\leq(\log T)^2$, the left side
has no pole near the real $s$-axis.  Its possible poles occur at
\begin{equation}\label{eq:Phi-remote-poles}
 s=n-\frac12-\gamma+it\quad(n\geq1),
\end{equation}
and the numerator gamma factors of $g$ add only
\begin{equation}\label{eq:g-remote-poles}
 s=-\frac12-\beta-it-2j,\qquad
 s=-\frac12-\gamma+it-2j\quad(j\geq0).
\end{equation}
All are at imaginary height comparable with $T$.
\end{lemma}

\begin{proof}
Duplication and reflection give
\[
 \chi(z)\Gamma(z)=\frac{(2\pi)^z}{2\cos(\pi z/2)}
\]
and
\[
 \frac{\Gamma(1-a_0-z)}{\Gamma(1-a_0)}
 =\frac{\Gamma(a_0)}{\Gamma(a_0+z)}
 \frac{\sin\pi a_0}{\sin\pi(a_0+z)}.
\]
Now use
\[
 \sin\pi(a_0+z)+\sin\pi a_0
 =2\sin\pi(a_0+z/2)\cos(\pi z/2).
\]
This proves \eqref{eq:gamma-collapse}, including cancellation of the apparent
odd-integer poles of $1/\cos(\pi z/2)$.  At nonpositive integral values of
$a_0+z$, the reciprocal gamma zero cancels the sine denominator.  The
remaining poles are \eqref{eq:Phi-remote-poles} and
\eqref{eq:g-remote-poles}.  If such a pole is crossed between fixed vertical
lines, the factor $e^{s^2}$ makes its residue $O_D(T^{-D})$.
\end{proof}

\begin{lemma}[Archimedean reflection]\label{lem:arch-reflection}
As meromorphic identities,
\begin{equation}\label{eq:Phi-at-zero}
               \Phi_{\lambda;\beta,\gamma}(0,t)
               =X_{\lambda,\gamma}(t),
\end{equation}
and
\begin{equation}\label{eq:Phi-reflection}
 X_{\beta,\gamma}(t)
 \Phi_{\lambda;-\gamma,-\beta}(s,t)
 =\Phi_{\lambda;\beta,\gamma}(-s,t).
\end{equation}
\end{lemma}

\begin{proof}
Substitute \eqref{eq:X-def}, \eqref{eq:chi-def}, and
\eqref{eq:Phi-def}.  Pair the full gamma factors by
\[
                  \Gamma(z)\Gamma(1-z)=\frac{\pi}{\sin\pi z}
\]
and the half-arguments by
\[
 \Gamma(z)\Gamma(z+\tfrac12)=2^{1-2z}\sqrt\pi\,\Gamma(2z).
\]
For \eqref{eq:Phi-at-zero}, the two terms in braces combine by reflection
and the remaining half-gamma factors form $X_{\lambda,\gamma}$.  For
\eqref{eq:Phi-reflection}, first combine the $\beta,\gamma$ factors with
$X_{\beta,\gamma}$ and then apply duplication.  The identities hold on a
pole-free open set and hence everywhere by meromorphic continuation.
\end{proof}

The endpoint-deleted direct and dual zero modes are therefore
\begin{align}
 \widetilde{\mathcal O}_{0,Y}^+(h,k)
 ={}&\int_{\R}w(t)\frac{1}{2\pi i}\int_{(c_u)}\Gamma(u)Y^u
 \frac{1}{2\pi i}\int_{(c_s)}\frac{\mathcal G(s,u)}s
 \Phi_{\lambda;\beta,\gamma}(s,t)\notag\\
 &\hspace{37mm}\times
 Z_{-\gamma-s,\beta+s,-\lambda}(h,k)\dd s\dd u\dd t,
 \label{eq:zero-plus-evaluated}\\
 \widetilde{\mathcal O}_{0,Y}^-(h,k)
 ={}&\int_{\R}w(t)\frac{1}{2\pi i}\int_{(c_u)}\Gamma(u)Y^u
 \frac{1}{2\pi i}\int_{(c_s)}\frac{\mathcal G(s,u)}s
 X_{\beta,\gamma}(t)
 \Phi_{\lambda;-\gamma,-\beta}(s,t)\notag\\
 &\hspace{37mm}\times
 Z_{\beta-s,-\gamma+s,-\lambda}(h,k)\dd s\dd u\dd t.
 \label{eq:zero-minus-evaluated}
\end{align}
The line $c_s>0$ is first chosen in the branchwise absolute-convergence
regions; elsewhere these formulas denote the continuation selected by the
full-kernel bridge of \Cref{sec:endpoint}.
